% Each feedback evaluates a stated step, not every preceding action.
\begin{tikzpicture}[
  x=1mm,y=1mm,font=\figurefont\fontsize{8}{10}\selectfont,
  box/.style={draw=BookRule,fill=BookPaper,text width=23mm,
    align=center,minimum height=10mm,inner sep=2mm,line width=.7pt},
  flow/.style={-{Stealth[length=2mm]},draw=BookMuted,line width=.8pt},
  signal/.style={flow,draw=BookAmber,dashed},
  lab/.style={fill=white,inner sep=.5mm,font=\figurefont\fontsize{7}{9}\selectfont}
]
  \node[box] (a0) at (15,23) {核对身份\\$a_0$};
  \node[box] (a1) at (58,23) {装载\\$a_1$};
  \node[box] (a2) at (101,23) {搬运（含选路）\\$a_2$};
  \node[box] (a3) at (144,23) {交付\\$a_3$};
  \draw[flow] (a0) -- (a1);
  \draw[flow] (a1) -- (a2);
  \draw[flow] (a2) -- (a3);
  \node[box,fill=FigureInput!10] (f0) at (15,-1) {人工标注\\读码动作};
  \node[box,fill=FigureInput!10] (f1) at (58,-1) {规则检查\\装载身份一致};
  \node[box,fill=FigureInput!10] (f2) at (101,-1) {模型步骤评价\\只评平稳程度};
  \node[box,fill=FigureInput!10] (f3) at (144,-1) {执行记录\\交付结果};
  \foreach \i in {0,1,2,3}{
    \draw[flow] (a\i) -- node[lab,right] {被评对象} (f\i);
    \node[box,fill=FigureModel!10] (u\i) at (\i*43+15,-26) {相应步骤的\\策略更新信号};
    \draw[signal] (f\i) -- node[lab,right] {反馈} (u\i);
  }
  \node[anchor=north,align=center] at (80,-38)
    {早期选路没有直接标签，仍靠结果回溯；终局身份与完好要求另行检验};
\end{tikzpicture}
